% !TeX root = ../quotients-and-descent.tex
\begin{thebibliography}{99}
\bibitem{AtiyahBott1968II}
M.~F.~Atiyah and R.~Bott, A Lefschetz fixed point formula for elliptic
complexes: II. Applications, \emph{Annals of Mathematics} (2)
\textbf{88} (1968), no.~3, 451--491; Theorem~4.12.
\url{https://webhomes.maths.ed.ac.uk/~v1ranick/papers/abott2.pdf}.

% PROVISIONAL: accompanying the separately reviewed ray-field claim.

\bibitem{CasalainaMartinKass2012}
S.~Casalaina-Martin and J.~L.~Kass,
\emph{A Riemann singularity theorem for integral curves},
Amer. J. Math. \textbf{134} (2012), no.~5, 1143--1165.
Introduction, classical smooth-curve statement; arXiv:0907.0212v2.
\url{https://arxiv.org/abs/0907.0212v2}

\bibitem{Goodson2024}
H. Goodson,
\emph{An exploration of degeneracy in abelian varieties of Fermat type},
arXiv:2211.03909v2, 1 August 2024.
The cited numbering refers to this version.
\url{https://arxiv.org/abs/2211.03909v2}

\bibitem{LiMantovanPries2021}
W.~Li, E.~Mantovan and R.~Pries,
\emph{Data for Shimura varieties intersecting the Torelli locus},
arXiv:2105.02286v1 (2021); Pacific J. Math. \textbf{343} (2026), 179--210.
Equation~(2.2), Proposition~6.11, Lemma~6.12 and Proposition~6.14
refer to the inspected arXiv version.
\url{https://arxiv.org/abs/2105.02286v1}

\bibitem{Masley1976}
J. M. Masley,
Solution of small class number problems for cyclotomic fields,
\emph{Compositio Mathematica} \textbf{33} (1976), 179--186,
\url{https://www.numdam.org/item/CM_1976__33_2_179_0.pdf}.

\bibitem{MilneModuli1972}
J. S. Milne,
Abelian varieties defined over their fields of moduli, I,
\emph{Bull. London Math. Soc.} \textbf{4} (1972), 370--372.
\url{https://www.jmilne.org/math/articles/1972c.pdf}

\bibitem{MilneModuli1974}
J. S. Milne,
Correction to ``Abelian varieties defined over their fields of moduli, I'',
\emph{Bull. London Math. Soc.} \textbf{6} (1974), 145--146,
including the addendum.

\bibitem{MilneAV2008}
J. S. Milne,
\emph{Abelian Varieties}, course notes, version 2.0, March 2008,
Chapter IV, Lemma 2.1, p.~135 (quotients by Galois-stable finite subgroups).
\url{https://www.jmilne.org/math/CourseNotes/AVc.pdf}

\bibitem{MilneShimura2017}
J.~S.~Milne, \emph{Introduction to Shimura Varieties},
revised September 16, 2017, Theorem~6.8 and Corollary~6.9,
\url{https://www.jmilne.org/math/xnotes/svi.pdf}.

\bibitem{MilneCFT2020}
J. S. Milne,
\emph{Class Field Theory}, version 4.03, 6 August 2020.
\url{https://www.jmilne.org/math/CourseNotes/CFT.pdf}

\bibitem{MilneCM2020}
J. S. Milne,
\emph{Complex Multiplication}, version 0.10, 14 July 2020,
Section 3.11 and Propositions 3.12--3.13 and 3.17, pp.~29--31.
\url{https://www.jmilne.org/math/CourseNotes/CM.pdf}

\bibitem{MilneJacobians2021}
J.~S.~Milne, \emph{Jacobian Varieties}, revised June12, 2021,
Theorem~12.1, p.~38; original chapter in \emph{Arithmetic Geometry}
(Springer, 1986), 167--212.
\url{https://www.jmilne.org/math/xnotes/JVs.pdf}

\bibitem{OpenAICM2026}
OpenAI math repository,
\emph{The rational Hodge conjecture for CM abelian varieties},
public manuscript dated 6 October 2026, introduction (main CM theorem)
and tensor-reduction section (auxiliary-factor contractions).
The comparison uses repository commit
\texttt{4eaee87548250b2682affda7d05d947c05554bae};
the general theorem and its formalization are not proof inputs here.
\url{https://github.com/openai/math/tree/4eaee87548250b2682affda7d05d947c05554bae/preprints/The-rational-Hodge-conjecture-for-CM-abelian-varieties-October-6-2026}

\bibitem{PolishchukTheta}
A.~Polishchuk, \emph{Abelian Varieties, Theta Functions and the Fourier
Transform}, Cambridge University Press, section~1.2, pp.~7--8.
Publisher excerpt: \url{https://assets.cambridge.org/97805218/08040/excerpt/9780521808040_excerpt.pdf}.

\bibitem{Saller2026}
K. Saller,
\emph{Split Weil structures on abelian sixfolds and the Hodge conjecture},
version 2, Zenodo record 23290280 (2026), Remarks~4.7--4.8.
\url{https://doi.org/10.5281/zenodo.23290280}

\bibitem{StacksDescent}
The Stacks Project Authors,
\emph{The Stacks Project}, Lemma 39.25.3 (tag 0CCJ),
and Section 35.35 (tag 02VV), finite locally free effective descent
and fpqc descent of morphisms.
\url{https://stacks.math.columbia.edu/tag/0CCJ};
\url{https://stacks.math.columbia.edu/tag/02VV}.

\bibitem{Zarhin2007}
Y.~G.~Zarhin, \emph{Cubic surfaces and cubic threefolds, jacobians
and intermediate jacobians}, arXiv:math/0610138v3 (2007), Theorem~1.1.
\url{https://math.nyu.edu/~tschinke/.manin/final/zarhin/zarhin.pdf}

\end{thebibliography}
